\chapter{Suivi des expériences et reproductibilité}\label{ch:ml:experiment-tracking-and-reproducibility}

Un régulateur demande à un établissement de montrer le modèle qui a servi au trading un matin donné, les données sur lesquelles il a été entraîné et chaque alternative
à laquelle il a été préféré. L'établissement peut nommer le modèle~; il faut trois semaines pour réentraîner quelque chose de proche, et il ne sait pas dire combien
d'alternatives il y en avait. Ce chapitre construit les enregistrements qui font de la réponse une requête~: un outil de suivi qui conserve chaque exécution
avec tout ce qui l'a déterminée, un registre qui indique quel modèle était en production à quel moment et d'où il provient, et une
\omterm{def:ml:experiment-tracking-and-reproducibility:registry}{piste d'audit} qui ne peut pas être modifiée sans que la modification apparaisse. Sur l'étude du chapitre, 200 exécutions de réglage d'une prévision par boosting, le
registre reproduit bit à bit le modèle de mars à partir de ses seules fiches~; si l'un quelconque des quatre champs manque, rien de proche
ne ressort. L'outil de suivi répond aussi à la question que le régulateur n'a pas posée~: en comptant les 200 essais, le ratio de Sharpe en validation du champion
vaut 2.22, et son ratio de Sharpe déflaté est de 0.15.

\section{Ce qui doit être enregistré}

\begin{definition}[Outil de suivi d'expériences, fiche d'exécution]\label{def:ml:experiment-tracking-and-reproducibility:tracker}
Un \emph{outil de suivi d'expériences}\index{outil de suivi d'expériences} conserve la trace de chaque exécution d'entraînement lancée par une équipe, y compris les
celles qui sont abandonnées. Une \emph{fiche d'exécution}\index{fiche d'exécution} rassemble
tout ce qui détermine le résultat de l'exécution (paramètres, graines, instantané de données, version du code, environnement) et
tout ce que l'exécution a produit (métriques, artefacts, journaux), avec l'heure et l'auteur.
\end{definition}

Le journal de recherche du livre~7 (chapitre~1) consignait les essais pour les compter~; l'outil les consigne pour les reconstruire. Dans le chapitre,
l'outil (le \cref{lst:ml:exp:run}) écrit chaque exécution sous la forme d'un enregistrement JSON dont l'identifiant est le hachage de son contenu, range le modèle ajusté
une seule fois sous le hachage de son contenu sérialisé (\texttt{firm.workflow} du livre~7), et enregistre l'exécution comme un essai dans
\texttt{firm.researchlog}. Les horodatages sont fournis en argument, jamais lus sur l'horloge, afin que le contenu de l'outil
lui-même soit reproductible.

L'étude~: une prévision par \omterm{def:ml:trees-and-boosting:boosting}{gradient boosting} (chapitre~5) des rendements du lendemain, sur un panel synthétique de 40 titres et huit variables sur
700 jours, entraînée sur les 400 premiers, évaluée par le ratio de Sharpe d'un portefeuille long-short quotidien sur les 150 suivants, et testée sur les
150 derniers. Deux cents exécutions tirent au hasard cinq paramètres et une graine. Leurs ratios de Sharpe en validation vont de $-3.16$ à 2.22,
avec une médiane de $-0.83$ (la \cref{fig:ml:exp:sharpes})~: le signal est faible, et la plupart des configurations ajustent du bruit.

\begin{omfigure}
\begin{tikzpicture}
\begin{axis}[width=11cm,height=5.8cm,ybar,bar width=12pt,xlabel={ratio de Sharpe en validation (annualisé)},ylabel={exécutions},
  tick label style={font=\small},label style={font=\small},xmin=-4,xmax=4,ymin=0,ymax=42,grid=major,grid style={gray!25},
  table/col sep=comma]
\addplot[fill=omProp!60,draw=none] table[x=x,y=runs]{figdata/ml/25-experiment-tracking-and-reproducibility/sharpes.csv};
\addplot[ycomb,omThm,very thick,mark=*,restrict expr to domain={\coordindex}{0:0}] table[x=x,y expr=40]{figdata/ml/25-experiment-tracking-and-reproducibility/marks.csv};
\addplot[ycomb,black,very thick,dashed,mark=none,restrict expr to domain={\coordindex}{1:1}] table[x=x,y expr=40]{figdata/ml/25-experiment-tracking-and-reproducibility/marks.csv};
\addplot[ycomb,omDef,very thick,mark=square*,restrict expr to domain={\coordindex}{2:2}] table[x=x,y expr=40]{figdata/ml/25-experiment-tracking-and-reproducibility/marks.csv};
\end{axis}
\end{tikzpicture}

\omcaption{Ratios de Sharpe en validation des 200 exécutions suivies~; cercle~: celui du champion (2.22)~; carré~: son ratio de Sharpe en test
(0.91)~; tirets~: le maximum espéré de 200 stratégies indépendantes sans compétence sur 150 jours (3.58). Données~:
\texttt{ml\_track.study}.}\label{fig:ml:exp:sharpes}
\end{omfigure}

Le ratio de Sharpe en test du champion vaut 0.91, soit moins de la moitié de son chiffre en validation. Son ratio de Sharpe probabiliste, c'est-à-dire la
probabilité que son ratio de Sharpe réel soit positif comme s'il était la seule exécution, vaut 0.96~; son ratio de Sharpe déflaté (livre~4,
chapitre~12), avec le nombre d'essais relevé dans l'outil, vaut 0.15. Le meilleur résultat espéré parmi 200 stratégies sans compétence sur 150 jours
est de 3.58, au-dessus du champion~: la correction est ici conservatrice, car les exécutions sont corrélées et forment moins d'essais indépendants
que 200, mais un résultat incapable de battre le meilleur de son propre bruit n'a pas mérité la mise en production sur la seule
preuve issue de la validation.

\section{Versionnage des données, du code et des modèles}

\begin{definition}[Versionnage des données, artefact de modèle]\label{def:ml:experiment-tracking-and-reproducibility:versions}
Le \emph{versionnage des données}\index{versionnage des données} identifie chaque jeu de données qu'une exécution lit au moyen d'un instantané immuable (un hachage du contenu ou une
version datée et figée), de sorte que les corrections ultérieures créent de nouvelles versions au lieu de modifier les anciennes. Un \emph{artefact
de modèle}\index{artefact de modèle} est le résultat stocké d'une exécution d'entraînement, les paramètres ajustés sous une forme portable, identifié par
son hachage de contenu.
\end{definition}

Les versions existent parce que les données et le code changent sous un modèle. Dans l'étude, le fournisseur de données révise 1~\% des rendements de mars
(l'instantané 2026-03 remplace 2026-01), et le code des variables modifie son écrêtage de trois écarts types à deux écarts types et demi
(v2 remplace v1). Le champion a été entraîné sur 2026-01 avec v1. La fiche le précise~; sans elle, une équipe utiliserait ce qu'elle a sous la main.

\section{Le registre et la piste d'audit}

\begin{definition}[Registre de modèles, lignage du modèle, piste d'audit]\label{def:ml:experiment-tracking-and-reproducibility:registry}
Un \emph{registre de modèles}\index{registre de modèles} recense les modèles de l'établissement par nom et version, avec leur étape (candidat,
production, retiré) et l'historique des changements d'étape. Le \emph{lignage du modèle}\index{lignage du modèle} relie une version enregistrée à
l'exécution, l'instantané des données et le code qui l'a produite. Une \emph{piste d'audit}\index{piste d'audit} est un journal d'actions
en ajout seul (enregistrements, promotions, dérogations) indiquant qui, quand et pourquoi, protégé contre les modifications silencieuses~; ici chaque entrée porte le hachage
de la précédente (le \cref{lst:ml:exp:audit}).
\end{definition}

\begin{dated}{2026-09}{Conservation des traces des modifications du trading algorithmique dans l'UE}\label{dat:ml:experiment-tracking-and-reproducibility:rts6}
Commission Delegated Regulation (EU) 2017/589 (les normes techniques de réglementation relatives aux exigences organisationnelles
du trading algorithmique, appelées RTS~6) oblige une société d'investissement à «~conserver les traces de toute modification substantielle apportée au logiciel
utilisé pour le trading algorithmique~», ce qui lui permet de déterminer quand une modification a été faite, qui l'a faite, qui l'a approuvée et quelle en est la nature
(article~5(7))~; les registres d'ordres des traders HFT sont conservés pendant cinq ans (article~28).
\end{dated}

Le registre répond à la première question du régulateur en fonction de la date (le \cref{lst:ml:exp:asof})~: la version en production le matin
du 10 mars se retrouve à partir de l'historique des étapes, et son lignage mène à la \omterm{def:ml:experiment-tracking-and-reproducibility:tracker}{fiche d'exécution}. La \omterm{def:ml:experiment-tracking-and-reproducibility:registry}{piste d'audit} répond à «~qui l'a promue,
et pourquoi~»~: la version 1 de \texttt{xs-gbdt}, promue le 2 mars comme champion de 200 exécutions.
Modifier un seul caractère de n'importe quelle ligne de la piste rompt la chaîne de hachage à cette ligne.

\section{Reproduire un modèle, et quand le bit à bit est trop exigeant}

À partir de la seule entrée du registre, le chapitre relance l'entraînement (le \cref{lst:ml:exp:repro}) et obtient un modèle dont la sérialisation a le
même hachage de contenu que l'artefact stocké~: bit à bit, le modèle qui a servi au trading. Il retire ensuite un champ de la fiche à la
fois et le remplace par ce qu'une équipe supposerait sans ce champ (le \cref{tab:ml:exp:ablation}). Chaque retrait casse la reproduction, et
aucun n'est anodin~: les ratios de Sharpe en test des modèles reproduits s'écartent de ceux de l'original de $-1.28$ avec les paramètres par défaut,
$-1.25$ avec les données révisées, $-0.44$ avec le code actuel des variables, et $+0.85$ avec la graine par défaut.

\begin{table}[htbp]
\centering\small
\begin{tabular}{@{}>{\raggedright\arraybackslash}p{3.6cm}>{\raggedright\arraybackslash}p{4.6cm}>{\centering\arraybackslash}p{2.4cm}>{\centering\arraybackslash}p{3.4cm}@{}}
\toprule
champ manquant & valeur supposée à la place & reproduit bit à bit & variation du ratio de Sharpe en test\\
\midrule
aucun (la fiche du registre) & & oui & 0\\
paramètres & valeurs par défaut de l'établissement & non & $-1.28$\\
graine & 1 & non & $+0.85$\\
instantané de données & la plus récente (2026-03) & non & $-1.25$\\
version du code & la version courante (v2) & non & $-0.44$\\
\bottomrule
\end{tabular}
\caption{Reproduction du modèle de production à partir de sa fiche, chaque champ étant tour à tour remplacé par une hypothèse. Données~:
\texttt{ml\_track.reproduction}.}\label{tab:ml:exp:ablation}
\end{table}

La ligne de la graine apporte un second enseignement~: une graine différente pour une même configuration fait varier le ratio de Sharpe du test de 0.85, soit autant que
l'écart entre deux configurations. Un résultat qui ne résiste pas à un changement de graine n'était pas un résultat, et rendre compte
de la fourchette sur les graines (chapitre~7) fait partie de ce qu'un outil de suivi permet.

Reproduire le modèle à l'identique, bit pour bit, est la bonne norme pour l'artefact propre d'un modèle, reconstruit sur le même logiciel et le même matériel~; les ajustements du chapitre
sont déterministes parce que LightGBM tourne sur un seul fil avec son option déterministe. D'une version de bibliothèque à l'autre, d'un nombre de fils
ou d'un matériel à l'autre (le \omterm{def:ml:training-infrastructure:mixed}{bfloat16} du chapitre~23), les sommes en virgule flottante changent d'ordre et les bits changent~; la norme est alors statistique~:
les mêmes métriques, à l'intérieur de la fourchette sur les graines, vérifiées en relançant le calcul. Ce qui doit toujours être exact, c'est la fiche~: les données, le code,
les paramètres et la graine qui ont produit l'artefact.

\begin{method}[Suivre la recherche qui trade]\label{met:ml:exp:method}
\begin{enumerate}
\item Enregistrer chaque exécution, y compris celles qui ont été abandonnées, avec ses paramètres, sa graine, son instantané de données, sa version du code et son environnement, et
stocker les artefacts par empreinte de contenu.
\item Prendre le nombre d'essais pour la déflation dans l'outil de suivi, et non de mémoire.
\item Enregistrer les modèles en production avec leur lignage et l'historique de leurs étapes~; ne les promouvoir et ne les retirer que par le registre, avec les motifs consignés dans
une \omterm{def:ml:experiment-tracking-and-reproducibility:registry}{piste d'audit} en ajout seul.
\item Tester la reproduction à partir du seul registre, régulièrement, et traiter tout échec comme un incident.
\end{enumerate}
\end{method}

\section{Tutoriel~: reproduire le modèle de mars}\label{tut:ml:experiment-tracking-and-reproducibility:tut}

\begin{tutorial}
\textbf{Objectif.} Suivre 200 exécutions de calibration, enregistrer et promouvoir le champion, le retrouver à partir de sa date, le reproduire bit pour bit, puis casser
la reproduction champ par champ. \textbf{État final~:} la \cref{fig:ml:exp:sharpes} et le \cref{tab:ml:exp:ablation}.
\begin{tutsteps}
\item \textbf{Une \omterm{def:ml:experiment-tracking-and-reproducibility:tracker}{fiche d'exécution}.}
\omcode{code/firm/exptrack/firm_exptrack.py}{73}{85}{Enregistrement d'une exécution~: artefact par empreinte de contenu, fiche par sa propre empreinte, essai dans le journal de recherche.}\label{lst:ml:exp:run}
\item \textbf{Une \omterm{def:ml:experiment-tracking-and-reproducibility:registry}{piste d'audit} en ajout seul.}
\omcode{code/firm/exptrack/firm_exptrack.py}{48}{64}{Entrées d'audit chaînées par hachage et leur vérification.}\label{lst:ml:exp:audit}
\item \textbf{Le modèle en production à une date donnée.}
\omcode{code/firm/exptrack/firm_exptrack.py}{134}{144}{Interrogation du registre à un instant donné.}\label{lst:ml:exp:asof}
\item \textbf{Reproduction.}
\omcode{code/firm/exptrack/firm_exptrack.py}{151}{155}{Nouvelle exécution à partir d'une fiche et comparaison des empreintes.}\label{lst:ml:exp:repro}
\item \textbf{Lancer} \texttt{ml\_track.summary()}, \texttt{reproduction()}, \texttt{null\_max()} et \texttt{fig\_track.py}
(environ 45 secondes sur un seul cœur).
\end{tutsteps}
\textbf{Pour aller plus loin.} Estimer le nombre effectif d'essais indépendants à partir de la corrélation des rendements de validation des exécutions
puis recalculer le ratio de Sharpe déflaté avec ce nombre.
\end{tutorial}

\section{Construction~: suivi des expériences}\label{bld:ml:experiment-tracking-and-reproducibility:exptrack}

\begin{build}
\bfield{Rôle} Chaque modèle que la société trade est traçable jusqu'à son exécution, ses données et son code, et peut être reconstruit à partir d'eux.
\bfield{Interface} \texttt{Tracker(root, log, family)} avec \texttt{run}, \texttt{get}, \texttt{artefact}, \texttt{runs},
\texttt{best}~; \texttt{Registry(root, audit)} avec \texttt{register}, \texttt{promote}, \texttt{current}, \texttt{as\_of},
\texttt{lineage}~; \texttt{AuditTrail(path)} avec \texttt{append} et \texttt{verify}~; \texttt{reproduce(record, train)}~; sur
\texttt{firm.workflow} et \texttt{firm.researchlog}.
\bfield{Règles} Les fiches ne sont écrites qu'une fois~; les artefacts sont adressés par leur contenu~; l'horloge n'entre jamais dans une fiche, sauf
sous la forme d'un horodatage fourni~; les promotions sont accompagnées de motifs.
\bfield{Tests d'acceptation} \texttt{code/firm/exptrack/tests/}~: des exécutions identiques partagent un artefact~; le registre renvoie la version
en production à une date passée~; promouvoir une nouvelle version retire l'ancienne~; une ligne d'audit falsifiée est détectée~; la reproduction réussit
à partir d'une fiche complète et échoue lorsqu'un champ change.
\bfield{Pour aller plus loin} Capture de l'environnement (versions des paquets, matériel) et un contrôle qui refuse de reproduire dans un environnement
différent~; instantanés de données sur \texttt{firm.pit}.
\end{build}

\begin{omsources}
\item J.~Pineau et coauteurs, ``Improving reproducibility in machine learning research (a report from the NeurIPS 2019
reproducibility program)'', arXiv:2003.12206, 2020.
\item O.~E.~Gundersen et S.~Kjensmo, ``State of the art: reproducibility in artificial intelligence'', \emph{AAAI}, 2018.
\item Commission Delegated Regulation (EU) 2017/589 du 19 juillet 2016 (RTS~6), \emph{Official Journal of the European Union}.
\end{omsources}

\section{Exercices}

\begin{exercise}[$\star$]\label{exo:ml:experiment-tracking-and-reproducibility:1}
Lister les cinq champs d'une \omterm{def:ml:experiment-tracking-and-reproducibility:tracker}{fiche d'exécution} qui déterminent son résultat, et une chose que chacun peut modifier à son insu.
\end{exercise}

\begin{exercise}[$\star$]\label{exo:ml:experiment-tracking-and-reproducibility:2}
Pourquoi l'identifiant d'une exécution est-il l'empreinte de sa fiche plutôt qu'un compteur ou un horodatage~?
\end{exercise}

\begin{exercise}[$\star$]\label{exo:ml:experiment-tracking-and-reproducibility:3}
Que signifie le ratio de Sharpe déflaté de 0.15, et pourquoi diffère-t-il du ratio de Sharpe probabiliste de 0.96~?
\end{exercise}

\begin{exercise}[$\star\star$]\label{exo:ml:experiment-tracking-and-reproducibility:4}
Pourquoi le maximum espéré de 200 ratios de Sharpe nuls est-il une référence trop sévère pour ces 200 exécutions~?
\end{exercise}

\begin{exercise}[$\star\star$]\label{exo:ml:experiment-tracking-and-reproducibility:5}
Comment une chaîne de hachage rend-elle une \omterm{def:ml:experiment-tracking-and-reproducibility:registry}{piste d'audit} détectable en cas de falsification, et que ne protège-t-elle pas~?
\end{exercise}

\begin{exercise}[$\star\star$]\label{exo:ml:experiment-tracking-and-reproducibility:6}
\emph{Trouver la faille.} «~Nous sauvegardons les poids de chaque modèle en production, donc nous pouvons toujours reproduire nos modèles.~»
\end{exercise}

\begin{exercise}[$\star\star\star$]\label{exo:ml:experiment-tracking-and-reproducibility:7}
\emph{Programmation.} Réentraîner la configuration du champion avec dix graines différentes et rendre compte de la fourchette des ratios de Sharpe du test.
Que dit-elle du champion~?
\end{exercise}

\begin{exercise}[$\star\star\star$]\label{exo:ml:experiment-tracking-and-reproducibility:8}
Concevoir ce que le registre et la \omterm{def:ml:experiment-tracking-and-reproducibility:registry}{piste d'audit} doivent contenir pour qu'une société respecte la consignation des modifications substantielles prévue par le RTS~6 pour un modèle
réentraîné chaque semaine.
\end{exercise}

\section{Problème~: reproduire le modèle de mars}

\begin{problem}[{Problème du week-end --- le modèle qui a tradé}]\label{pb:ml:experiment-tracking-and-reproducibility:1}
Les 200 exécutions du chapitre, le registre et la \omterm{def:ml:experiment-tracking-and-reproducibility:registry}{piste d'audit}.

\medskip\noindent\textbf{Partie I --- L'étude.}
\begin{enumerate}
\item Que contient la fiche de chaque exécution~?
\item Quelle est la distribution des ratios de Sharpe de validation, et quels sont les chiffres du champion~?
\item Quels sont ses ratios de Sharpe probabiliste et déflaté~?
\item Quel est le meilleur espéré de 200 stratégies nulles, et pourquoi est-il supérieur au champion~?
\end{enumerate}

\medskip\noindent\textbf{Partie II --- Versions et registre.}
\begin{enumerate}[resume]
\item Qu'est-ce qui a changé dans les données et le code après l'entraînement du champion~?
\item Comment retrouve-t-on le modèle en production le 10 mars~?
\item Que consigne la \omterm{def:ml:experiment-tracking-and-reproducibility:registry}{piste d'audit}, et comment détecte-t-on une altération~?
\item Que impose le RTS~6 à la société en matière d'enregistrement~?
\end{enumerate}

\medskip\noindent\textbf{Partie III --- Reproduction.}
\begin{enumerate}[resume]
\item Que se passe-t-il lorsque le modèle est reconstruit à partir de sa fiche~?
\item Que se passe-t-il lorsque chaque champ est absent~?
\item Que montre la ligne de la graine~?
\item Quand la reproduction bit à bit est-elle la mauvaise norme~?
\end{enumerate}

\medskip\noindent\textbf{Partie IV --- Le verdict.}
\begin{enumerate}[resume]
\item Énoncer le \emph{résultat nommé}~: le nombre d'essais derrière le champion et son ratio de Sharpe déflaté, ainsi que le champ dont
l'absence a fait échouer la reproduction.
\item Auriez-vous promu le champion~? Qu'auriez-vous demandé~?
\item Comment répondriez-vous à la demande du régulateur avec ces outils, et combien de temps cela prendrait-il~?
\item Que faudrait-il ajouter au suivi pour satisfaire une règle de conservation de cinq ans~?
\item Comment le suivi doit-il compter les essais qui ont été menés mais jamais enregistrés~?
\item Que doit contenir le registre pour un modèle réentraîné automatiquement chaque semaine~?
\item Qui doit être autorisé à promouvoir un modèle~?
\item En une phrase~: que faut-il pour reproduire un modèle~?
\end{enumerate}
\end{problem}

\section{Questions d'entretien}

\begin{interviewq}[$\star$ \iqroles{ingénieur ML}]\label{iq:ml:experiment-tracking-and-reproducibility:1}
Que enregistre un \omterm{def:ml:experiment-tracking-and-reproducibility:tracker}{outil de suivi d'expériences}, et pourquoi enregistrer les exécutions échouées~?
\end{interviewq}

\begin{interviewq}[$\star\star$ \iqroles{ingénieur ML, développeur}]\label{iq:ml:experiment-tracking-and-reproducibility:2}
Comment rendez-vous un entraînement reproductible bit pour bit~?
\end{interviewq}

\begin{interviewq}[$\star\star$ \iqroles{chercheur}]\label{iq:ml:experiment-tracking-and-reproducibility:3}
Comment le nombre d'essais modifie-t-il la signification du ratio de Sharpe d'un backtest~?
\end{interviewq}

\begin{interviewq}[$\star\star$ \iqroles{ingénieur ML}]\label{iq:ml:experiment-tracking-and-reproducibility:4}
Concevez un \omterm{def:ml:experiment-tracking-and-reproducibility:registry}{registre de modèles} pour une société de trading. Quelles étapes et quels champs doit-il comporter~?
\end{interviewq}

\begin{interviewq}[$\star\star$ \iqroles{ingénieur ML, développeur}]\label{iq:ml:experiment-tracking-and-reproducibility:5}
Pourquoi versionner les données, et comment versionneriez-vous un jeu de données corrigé a posteriori~?
\end{interviewq}

\begin{interviewq}[$\star\star\star$ \iqroles{ingénieur ML}]\label{iq:ml:experiment-tracking-and-reproducibility:6}
Votre modèle en production ne peut pas être reproduit. Décrivez l'enquête pas à pas.
\end{interviewq}
